%% ============================================================
\section{بخش سوم: پردازش تصویر، ایمیل و \lr{MQTT}}
%% ============================================================

\subsection{تشخیص انسان}

مدل مورد استفاده \lr{MobileNet-SSD} آموزش‌دیده روی \lr{PASCAL VOC} است که از طریق \lr{OpenCV DNN} با \lr{backend} و \lr{target} روی \lr{CPU} اجرا می‌شود. کلاس شماره‌ی ۱۵ در این مجموعه «انسان» است. روی هر فریم خروجی، کادر دور افراد، شمارنده، شماره‌ی دانشجویی، تاریخ، ساعت زنده‌ی سیستم و \lr{FPS} اندازه‌گیری‌شده نوشته می‌شود.

\subsubsection*{دروازه‌ی حرکت}

اجرای شبکه روی هر فریم روی این سخت‌افزار ممکن نیست. به همین دلیل پیش از هر استنتاج، یک تصویر بندانگشتی \lr{64x48} خاکستری از فریم ساخته می‌شود و با بندانگشتی فریم قبلی مقایسه می‌گردد؛ اگر میانگین تفاضل مطلق از آستانه کمتر باشد، صحنه بدون تغییر تلقی شده و نتیجه‌ی قبلی بازاستفاده می‌شود. در عمل این دروازه حدود ۹۰ درصد فریم‌ها را حذف می‌کند.

برای آنکه یک صحنه‌ی کاملاً ساکن باعث نشود سامانه برای همیشه به نتیجه‌ی قدیمی تکیه کند، یک \lr{heartbeat} هم وجود دارد که صرف‌نظر از حرکت، هر چند ثانیه یک استنتاج اجباری انجام می‌دهد.

\subsubsection*{مشکل مهمی که کشف و برطرف شد: کندی و ناپایداری استریم}

در آزمایش‌های اولیه، استریم رفتار عجیبی داشت: وقتی اتاق خالی بود حدود ۵ فریم بر ثانیه کار می‌کرد، اما \emph{به‌محض آنکه فردی وارد کادر می‌شد} به حدود ۱٫۵ فریم بر ثانیه سقوط می‌کرد. یعنی سامانه دقیقاً در لحظه‌ای که باید کار می‌کرد بدترین عملکرد را داشت.

علت پس از بررسی مشخص شد: فراخوانی \lr{net.forward()} به‌صورت \textbf{هم‌زمان (\lr{synchronous})} در همان حلقه‌ای بود که فریم‌ها را منتشر می‌کرد. هر استنتاج حدود ۱٫۲۵ ثانیه طول می‌کشد و در تمام این مدت هیچ فریمی منتشر نمی‌شد. وقتی اتاق خالی است دروازه‌ی حرکت ۹۲ درصد فریم‌ها را رد می‌کند و تقریباً استنتاجی رخ نمی‌دهد؛ اما وقتی فردی وارد می‌شود، \emph{حرکت می‌کند}، دروازه باز می‌شود و هر استنتاج حلقه را ۱٫۲۵ ثانیه قفل می‌کند.

راه‌حل، انتقال استنتاج به یک \textbf{نخ کارگر جداگانه} با صف تک‌خانه‌ای بود. تنها یک فریم در هر لحظه در حال پردازش است و همیشه تازه‌ترین فریم؛ صف عمیق‌تر فقط تأخیر انباشته می‌کرد، چون فریمی که مدتی منتظر مانده در برابر فریمی که همین الان رسیده بی‌ارزش است.

\begin{center}
\begin{tabular}{lrr}
\hline
 & پیش از اصلاح & پس از اصلاح \\
\hline
\lr{FPS} با حضور فرد & ۱٫۵۴ & ۵٫۹ تا ۷٫۵ \\
زمان استنتاج & ۱۲۷۹ میلی‌ثانیه & ۱۲۷۹ میلی‌ثانیه \\
درصد فریم‌های حذف‌شده & ۹۲٪ & ۹۰٪ \\
\hline
\end{tabular}
\end{center}

نکته‌ی کلیدی این است که \textbf{زمان استنتاج تغییری نکرده}؛ تشخیص سریع‌تر نشده، بلکه مسیر انتشار تصویر دیگر منتظر آن نمی‌ماند. بهبود حدود چهار برابر، دقیقاً در همان حالتی که پیش‌تر بدترین عملکرد را داشت.

معامله‌ای که این تغییر تحمیل می‌کند و باید صادقانه ذکر شود: کادرها از آخرین استنتاج \emph{کامل‌شده} رسم می‌شوند، پس روی سوژه‌ی متحرک کادر تا حدود ۱٫۲۵ ثانیه عقب‌تر از بدن حرکت می‌کند. البته این تأخیر پیش از این هم وجود داشت -- نتایج میان استنتاج‌ها بازاستفاده می‌شدند -- فقط حالا روی تصویر روان دیده می‌شود به‌جای آنکه در یک اسلایدشو پنهان بماند.

\subsection{ارسال ایمیل}

ارسال ایمیل کاملاً در کد \lr{C} و با \lr{libcurl} روی \lr{smtps://smtp.gmail.com:465} انجام می‌شود. پیام شامل تعداد افراد، برچسب زمانی، دمای لحظه‌ای پردازنده و تصویر پیوست‌شده‌ی لحظه‌ی تشخیص است.

\subsubsection*{سازوکار \lr{debounce}}

صورت پروژه حداکثر یک ایمیل در هر ۳۰ ثانیه را الزام کرده است. پیاده‌سازی بر پایه‌ی زمان آخرین ارسال \emph{به تفکیک نوع رویداد} است، نه یک زمان‌سنج سراسری. دلیل این تفکیک آن است که یک هشدار «دستکاری دوربین» نباید به این خاطر که چند ثانیه پیش یک ایمیل تشخیص عادی ارسال شده، بلعیده شود؛ این دو رویداد فوریت یکسانی ندارند.

در حالت نگهبانی، پنجره‌ی \lr{debounce} به ۳ ثانیه کاهش می‌یابد. منطق این است که در حالت عادی هدف \lr{debounce} جلوگیری از سیل ایمیل برای یک نفر که در اتاق راه می‌رود است، اما در حالت نگهبانی هر تشخیص یک رویداد امنیتی است و تأخیر ۳۰ ثانیه‌ای معنا ندارد؛ با این حال یک کف ضدسیل لازم است.

در گزارش \lr{journal} هر دو حالت دیده می‌شود:

\begin{latin}
\begin{verbatim}
mailer: "[Guardian 402170516] thermal throttling at 67.1 C" delivered in 3.3s
mailer: '[Guardian 402170516] thermal throttling at 67.1 C' suppressed
        (same kind sent 10s ago)
\end{verbatim}
\end{latin}

\subsection{\lr{MQTT}}

\lr{Broker} از نوع \lr{Mosquitto} روی کامپیوتر شخصی اجرا می‌شود و کلاینت برد به زبان \lr{C} با \lr{libmosquitto} نوشته شده است. ساختار موضوع‌ها مطابق الزام صورت پروژه است:

\begin{center}
\begin{tabular}{ll}
\hline
موضوع & محتوا \\
\hline
\lr{telemetry/402170516/home} & دما، حافظه، بار پردازنده، \lr{timestamp} \\
\lr{persons/402170516/home} & تعداد افراد، دما، \lr{timestamp} \\
\lr{status/402170516/home} & وضعیت \lr{online}/\lr{offline} (\lr{LWT}، \lr{retained}) \\
\lr{alarm/402170516/home} & هشدار حالت نگهبانی \\
\hline
\end{tabular}
\end{center}

پیام‌ها با \lr{QoS = 1} ارسال می‌شوند و \lr{LWT} روی موضوع \lr{status} با پرچم \lr{retained} تنظیم شده است.

\subsubsection*{مشکلی که در حین کار پیش آمد: آدرس \lr{Broker}}

آدرس \lr{Broker} ابتدا به‌صورت یک \lr{IP} ثابت در پیکربندی نوشته شده بود. اما لپ‌تاپ آدرس خود را از \lr{DHCP} می‌گیرد و پس از اتصال مجدد به شبکه، آدرسش تغییر کرد. نتیجه این بود که کلاینت در هر بار راه‌اندازی حدود ۲۵ تا ۴۰ ثانیه تلاش ناموفق می‌کرد تا سرانجام به آدرس پشتیبان برگردد و در این مدت تمام پیام‌های \lr{telemetry} از دست می‌رفت.

راه‌حل، استفاده از نام \lr{Bonjour} لپ‌تاپ به‌جای \lr{IP} بود. بسته‌ی \lr{libnss-mdns} روی برد نصب است و \lr{avahi-daemon} این نام را برای کاربر \lr{guardian} \lr{resolve} می‌کند، بنابراین نام، لپ‌تاپ را در هر آدرسی که بگیرد دنبال می‌کند. زمان اتصال از ۲۵ تا ۴۰ ثانیه به \textbf{۹ میلی‌ثانیه} کاهش یافت.

\subsection{نتایج آزمایش‌های بخش سوم}

\subsubsection*{آزمایش ۳--۴: \lr{LWT} و قطع \lr{Broker}}

نکته‌ی مهم مفهومی: پیام \lr{LWT} را \emph{خود \lr{Broker}} از طرف کلاینتی که بدون خداحافظی ناپدید شده منتشر می‌کند. بنابراین خاموش‌کردن \lr{Broker} به‌تنهایی هرگز نمی‌تواند \lr{LWT} را نشان دهد، چون دیگر کسی باقی نمانده که آن را منتشر کند. به همین دلیل آزمایش در دو بخش انجام شد:

\begin{enumerate}
	\item کشتن دیمن برد با \lr{SIGKILL} در حالی که \lr{Broker} فعال است. پیام \lr{offline} پس از یک بازه‌ی \lr{keepalive} (۲۰ ثانیه) روی موضوع \lr{status} دریافت شد.
	\item خاموش‌کردن \lr{Broker} به مدت ۳ دقیقه. برد \lr{mqtt\_connected=false} گزارش می‌دهد، با \lr{backoff} تلاش مجدد می‌کند، و \textbf{وب‌سرور و \lr{API} و آشکارساز بدون وقفه به کار ادامه می‌دهند}. پس از بازگشت \lr{Broker}، کلاینت خودبه‌خود و بدون دخالت متصل شد.
\end{enumerate}

\subsubsection*{آزمایش ۳--۶: ورود غیرمجاز \lr{MQTT}}

هر چهار تلاش رد شد:

\begin{center}
\begin{tabular}{ll}
\hline
تلاش & نتیجه \\
\hline
اتصال ناشناس (بدون نام کاربری) & \lr{Connection Refused: not authorised} \\
کاربر معتبر با رمز اشتباه & \lr{Connection Refused: not authorised} \\
کاربر ناشناخته & \lr{Connection Refused: not authorised} \\
انتشار توسط کاربر فقط‌خواندنی & اتصال برقرار، پیام توسط \lr{ACL} دور ریخته شد \\
\hline
\end{tabular}
\end{center}

مورد چهارم نیازمند روش متفاوتی برای اثبات بود. وقتی \lr{ACL} یک انتشار را رد می‌کند، \lr{Broker} اتصال را می‌پذیرد، پیام را می‌گیرد و بی‌صدا دور می‌ریزد؛ \lr{mosquitto\_pub} هم در هر دو حالت کد خروج صفر برمی‌گرداند. بنابراین آزمون به‌صورت \emph{رفتاری} طراحی شد: یک پیام با نشانگر یکتا توسط کاربر فقط‌خواندنی منتشر شد در حالی که یک مشترک مجاز گوش می‌داد، و شرط قبولی آن بود که آن نشانگر هرگز بازنگردد در حالی که پیام‌های واقعی برد در همان مدت دریافت می‌شدند.

در تمام مدت این آزمایش، اتصال کلاینت مجاز برد به‌عنوان شاهد بررسی شد تا مطمئن شویم رد شدن تلاش‌ها به دلیل رد اعتبارنامه است نه به دلیل در دسترس نبودن \lr{Broker}.

\subsubsection*{آزمایش ۳--۷: ورود غیرمجاز \lr{SSH}}

سیاست فعال: \lr{PasswordAuthentication no}، \lr{PermitRootLogin no}، \lr{AllowUsers mahan}، \lr{MaxAuthTries 3}.

\begin{center}
\begin{tabular}{ll}
\hline
تلاش & نتیجه \\
\hline
ورود با رمز عبور & \lr{Permission denied (publickey).} \\
ورود با کاربر \lr{root} & \lr{Permission denied (publickey).} \\
کاربر ناشناخته & \lr{Permission denied (publickey).} \\
\hline
\end{tabular}
\end{center}

نکته آنکه \lr{sshd} اصلاً درخواست رمز عبور نمایش نمی‌دهد، پس چیزی برای \lr{brute-force} وجود ندارد. کلید مجاز در همان اجرا با موفقیت وارد شد تا ثابت شود ردها ناشی از سیاست است نه از دسترس‌ناپذیری برد.

%% ============================================================
\section{بخش چهارم: قابلیت‌های تکمیلی}
%% ============================================================

\subsection{حالت نگهبانی}

از طریق \lr{API} و دکمه‌ای روی صفحه‌ی \lr{HTML} فعال و غیرفعال می‌شود. در حالت فعال، هر تشخیص انسان باعث ارسال فوری ایمیل با عکس و انتشار پیام اضطراری روی \lr{alarm/402170516/home} می‌شود.

این دستور، برخلاف \lr{endpoint}های خواندنی، نیازمند \lr{Bearer token} است. توکن در فایلی با دسترسی \lr{root} نگهداری می‌شود و هرگز در کد ظاهر نمی‌شود.

\subsection{جعبه‌ی سیاه}

تاریخچه‌ی تشخیص‌ها در \lr{SQLite} با بافر حلقوی ذخیره می‌شود. کران‌داری اندازه با یک \lr{trigger} از نوع \lr{AFTER INSERT} اعمال می‌شود که هر چیزی قدیمی‌تر از \lr{db\_ring\_size} رکورد آخر را حذف می‌کند.

نکته‌ی طراحی مهم: شمارنده‌های تجمعی در جدول \emph{جداگانه‌ای} به نام \lr{stats} نگهداری می‌شوند، دقیقاً به این دلیل که از بافر حلقوی جان سالم به در ببرند. اگر شمارش کل در همان جدول \lr{detections} بود، پاسخ پرسش «امروز چند نفر دیده شدند» هر ۱۰۰۰ تشخیص بی‌صدا صفر می‌شد.

نتیجه‌ی آزمایش ۴--۲: ۶۰۷ رکورد ذخیره‌شده، درون کران ۱۰۰۰؛ \lr{endpoint} تاریخچه رکوردها را بازمی‌گرداند.

\subsection{\lr{Watchdog} نرم‌افزاری}

اگر بیش از ۳۰ ثانیه فریم تازه‌ای نرسد، سامانه ایمیل هشدار «دستکاری دوربین» می‌فرستد و سرویس آشکارساز را ری‌استارت می‌کند.

جزئیات ظریفی که در پیاده‌سازی اهمیت دارد: \lr{watchdog} سن آخرین فریمی را می‌سنجد که \emph{ورودی واقعی دوربین} داشته است، نه سن آخرین فریم منتشرشده. تمایز این دو حیاتی است، چون آشکارساز در نبود سیگنال یک تصویر جایگزین با نرخ ۱ هرتز منتشر می‌کند تا داشبورد یخ نزند؛ یک بررسی ساده‌لوحانه‌ی تازگی، همین فریم‌های مصنوعی را برای همیشه تازه می‌دید و هرگز فعال نمی‌شد.

نتیجه‌ی آزمایش ۴--۳:

\begin{center}
\begin{tabular}{ll}
\hline
\lr{PID} آشکارساز پیش از رویداد & ۳۴۹۵۱ \\
\lr{PID} آشکارساز پس از رویداد & ۳۵۶۷۹ \\
ایمیل هشدار & \lr{"CAMERA TAMPERING suspected"} در ۲٫۸ ثانیه \\
وضعیت دیمن اصلی در تمام مدت & \lr{active}، پاسخ \lr{HTTP 200} \\
\hline
\end{tabular}
\end{center}

معیار سنجش بازیابی، تغییر \lr{PID} است نه شمارنده‌ی \lr{NRestarts} سیستم‌دی. دلیلش این است که \lr{watchdog} بازیابی را با درخواست \lr{restart} از \lr{systemd} از طریق \lr{polkit} انجام می‌دهد، و \lr{NRestarts} تنها ری‌استارت‌هایی را می‌شمارد که خود \lr{systemd} پس از خرابی و بر اساس \lr{Restart=} آغاز کرده باشد -- یعنی در یک بازیابی کاملاً موفق \lr{watchdog} روی صفر می‌ماند و گزارش‌کردن آن گمراه‌کننده است.

سطر آخر جدول همان چیزی است که اصلاح وابستگی یونیت‌ها را اثبات می‌کند: دیمن اصلی در حالی فعال ماند که آشکارساز زیر پایش ری‌استارت می‌شد.

\subsection{مدیریت حرارتی تطبیقی}

آستانه‌های بالا و پایین جداگانه‌اند: ۷۰ درجه برای بالا رفتن سطح و ۶۵ درجه برای بازگشت. استفاده از دو آستانه به‌جای یکی عمدی است؛ با یک آستانه‌ی واحد سامانه نوسان می‌کرد: به‌محض کاهش نرخ فریم دما زیر خط می‌آمد، نرخ کامل بازمی‌گشت، دوباره گرم می‌شد و این چرخه دقیقه‌ای چند بار رزولوشن را تغییر می‌داد. فاصله‌ی ۵ درجه‌ای باعث می‌شود هر پله متعهد بماند.

\subsubsection*{آزمایش ۴--۴ و محدودیتی که آشکار کرد}

اجرای \lr{stress-ng} روی هر چهار هسته به مدت سه دقیقه، دما را تنها به \textbf{۶۸٫۷۶ درجه} رساند -- زیر آستانه‌ی ۷۰ درجه. یعنی این آستانه روی این برد \emph{از نظر فیزیکی دست‌نیافتنی} است، چون منبع تغذیه‌ی ناکافی پیش از آنکه حد حرارتی خود تراشه اهمیت پیدا کند، فرکانس را محدود می‌کند.

یک حلقه‌ی کنترل را نمی‌توان با انتظار برای نقطه‌ی تنظیمی که سخت‌افزار توان رسیدن به آن را ندارد راستی‌آزمایی کرد. بنابراین نقطه‌ی تنظیم به سخت‌افزار نزدیک شد: آستانه‌ها موقتاً به ۶۶٫۸ و ۶۲٫۸ درجه کاهش یافتند، بار دوباره اعمال شد و رفتار حاکم مشاهده و سپس مقادیر اصلی بازگردانده شدند:

\begin{latin}
\begin{verbatim}
thermal: level 2 (step 2) at 67.1 C -> 3 fps, 320 px capture, 224 px network
thermal: level 3 (step 3) at 67.1 C -> 2 fps, 320 px capture, 192 px network
mailer: "[Guardian 402170516] thermal throttling at 67.1 C" delivered in 3.3s
\end{verbatim}
\end{latin}

آنچه این آزمایش نشان می‌دهد \emph{سازوکار} است: تشخیص، کاهش پله‌ای نرخ فریم و رزولوشن، اطلاع‌رسانی و بازگشت. آنچه نشان \emph{نمی‌دهد} رسیدن برد به ۷۰ درجه است، که با این منبع تغذیه ممکن نیست.

%% ============================================================
\section{وضعیت منبع تغذیه: یک محدودیت اندازه‌گیری‌شده}
%% ============================================================

در جریان آزمایش‌ها مشخص شد که برد دچار افت ولتاژ است. در شش ساعت نخست کارکرد، هسته‌ی لینوکس \textbf{۳۵۱ رویداد \lr{Undervoltage detected!}} ثبت کرد و \lr{vcgencmd get\_throttled} مقدار \lr{0x50005} را برمی‌گرداند که بیت‌های آن یعنی: افت ولتاژ هم‌اکنون، محدودسازی فعال هم‌اکنون، و سابقه‌ی هر دو از زمان بوت.

دمای پردازنده در آن لحظه حدود ۶۲ درجه بود، یعنی بسیار پایین‌تر از حد نرم ۸۰ درجه؛ پس این محدودسازی \emph{حرارتی نیست}، بلکه ناشی از ولتاژ است. نمونه‌برداری از فرکانس نشان داد هسته‌ها عمدتاً روی ۱۴۰۰ مگاهرتز کار می‌کنند اما گاهی به ۱۱۰۰ تا ۱۲۰۰ مگاهرتز افت می‌کنند.

اثر این پدیده بر نتایج به‌صورت شفاف در گزارش ذکر شده است:

\begin{itemize}
	\item زمان استنتاج حدود ۱۵ تا ۲۰ درصد در بازه‌های افت بدتر می‌شود. علت اصلی کندی استنتاج \emph{این نیست} -- مدل \lr{MobileNet-SSD} با ورودی ۳۰۰ پیکسل و اعداد اعشاری روی هسته‌ی \lr{Cortex-A53} ذاتاً همین‌قدر زمان می‌برد.
	\item آزمایش ۴--۴ نتوانست آستانه‌ی طراحی‌شده را فعال کند و ناچار به روش نقطه‌ی تنظیم کاهش‌یافته شد.
	\item دماهای گزارش‌شده در آزمایش ۲--۱ سقفی پایین‌تر از سیلیکون بدون محدودیت دارند؛ ترتیب \emph{نسبی} سه حالت همچنان معتبر است.
\end{itemize}

راه‌حل، یک منبع تغذیه‌ی استاندارد ۵٫۱ ولت با جریان حداقل ۲٫۵ آمپر و کابل کوتاه و ضخیم است. تمام اندازه‌گیری‌ها با علم به این محدودیت انجام و گزارش شده‌اند.

%% ============================================================
\section{الزامات امنیتی}
%% ============================================================

\subsection{ورود \lr{SSH}}

ورود تنها با کلید عمومی ممکن است و ورود \lr{root} غیرفعال است. اسکریپت \lr{harden\_ssh.sh} این تغییر را با سه لایه‌ی ایمنی اعمال می‌کند: اثبات کارکرد احراز هویت با کلید، اعتبارسنجی پیکربندی با \lr{sshd -t}، و مسلح‌کردن یک زمان‌سنج بازگشت خودکار ده‌دقیقه‌ای که تنها با تأیید صریح خلع سلاح می‌شود.

در جریان اجرای این اسکریپت دو ایراد واقعی در خودش کشف و برطرف شد:

\begin{enumerate}
	\item \textbf{بررسی ایمنی نادرست بود.} اسکریپت برای اطمینان از کارکرد کلید، یک ورود \lr{loopback} از برد به خودش انجام می‌داد. این کار مستلزم نگهداری یک کلید \emph{خصوصی} روی همان دستگاهی است که قرار است ایمن شود، و در ضمن مسیری را می‌آزماید که اهمیتی ندارد. بررسی به خواندن گزارش خود \lr{sshd} تغییر یافت: وجود سطر \lr{Accepted publickey} برای کاربر مجاز، اثبات مستقیم همان مسیری است که باید پس از قطع رمز عبور کار کند.

	\item \textbf{یک باگ \lr{shell} که نتیجه را وارونه می‌کرد.} الگوی \lr{journalctl | grep -q} زیر \lr{set -o pipefail} نادرست است: \lr{grep -q} در نخستین تطبیق خارج می‌شود، \lr{journalctl} سیگنال \lr{SIGPIPE} می‌گیرد و کد خروج ۱۴۱ برمی‌گرداند، و \lr{pipefail} کل لوله را ناموفق اعلام می‌کند. یعنی \emph{تطبیق موفق} به‌صورت «تطبیقی یافت نشد» تفسیر می‌شد. همین الگو در مسیر خلع سلاح هم بود و باعث می‌شد زمان‌سنج بازگشت خودکار بی‌صدا متوقف نشود.
\end{enumerate}

\subsection{\lr{MQTT}}

پیکربندی \lr{Mosquitto} شامل \lr{allow\_anonymous false}، فایل رمز عبور و فایل \lr{ACL} است. کاربر \lr{guardian} تنها اجازه‌ی \lr{write} روی موضوعات پروژه و کاربر \lr{viewer} تنها اجازه‌ی \lr{read} دارد.

\subsection{نبود اعتبارنامه در کد}

هیچ رمز عبور یا کلیدی در کد یا در مخزن وجود ندارد. اعتبارنامه‌ها در \lr{/etc/guardian/secrets.env} با مالکیت \lr{root:guardian} و دسترسی \lr{0640} نگهداری می‌شوند و از طریق \lr{EnvironmentFile} به سرویس تزریق می‌گردند.

اسکریپت \lr{check\_secrets.sh} به‌عنوان سند این ادعا نوشته شده و سه چیز را بررسی می‌کند: نبود اعتبارنامه در درخت کد، نبود آن در \emph{باینری کامپایل‌شده} (که یک \lr{grep} روی کد منبع از قلم می‌انداخت)، و درستی مالکیت و دسترسی فایل اسرار.

%% ============================================================
\section{جمع‌بندی و مشکلات پیش‌آمده}
%% ============================================================

مهم‌ترین مسائلی که در طول پیاده‌سازی با آن‌ها روبه‌رو شدم و نحوه‌ی حل آن‌ها:

\begin{enumerate}
	\item \textbf{زنجیره‌ی وابستگی مرگبار در \lr{systemd}.} وابستگی \lr{Requires} به سوکتی که \lr{PartOf} آشکارساز بود، باعث می‌شد ری‌استارت آشکارساز کل سامانه را از بین ببرد. با تنزل به \lr{After} حل شد. این ایراد در آزمایش ۲--۱ خود را نشان داد و اگر کشف نمی‌شد، آزمایش ۴--۳ سامانه را برای همیشه از کار می‌انداخت.

	\item \textbf{استنتاج هم‌زمان در مسیر انتشار تصویر.} علت سقوط \lr{FPS} از ۵ به ۱٫۵ هنگام حضور فرد. با انتقال به نخ کارگر و صف تک‌خانه‌ای، حدود چهار برابر بهبود یافت.

	\item \textbf{آدرس \lr{IP} ثابت برای \lr{Broker}.} با تغییر شبکه‌ی لپ‌تاپ بی‌اعتبار شد. با نام \lr{Bonjour} جایگزین شد و زمان اتصال از ۲۵ ثانیه به ۹ میلی‌ثانیه رسید.

	\item \textbf{خطای \lr{ffmpeg} در گرفتن تصویر وب‌کم.} ورودی \lr{avfoundation} تنها زمانی یک حالت را می‌پذیرد که نرخ فریم درخواستی برابر \emph{بیشینه}‌ی آن حالت باشد؛ بنابراین درخواست ۱۵ فریم بر ثانیه از دوربینی که \lr{[15 30]} را اعلام می‌کند رد می‌شود، هرچند ۱۵ واقعاً پشتیبانی می‌شود. علاوه بر این، \lr{pixel\_format} درخواستی روی دوربین‌های \lr{Apple Silicon} موجود نبود. هر دو خطا با پیام مبهم \lr{Input/output error} ظاهر می‌شدند. راه‌حل: اسکریپت اکنون در آغاز از خود دستگاه پرس‌وجو می‌کند و حالت‌ها و \lr{pixel format}های پشتیبانی‌شده را استخراج و انتخاب می‌کند.

	\item \textbf{خرابی داده در اسکریپت‌های اندازه‌گیری.} فرمان \lr{vcgencmd} بدون \lr{root} پیامی \emph{دوخطی} روی \lr{stderr} چاپ می‌کند و آن خط جدید هر سطر \lr{CSV} را نصف می‌کرد. همچنین محاسبه‌ی شیب حداقل مربعات با استفاده از زمان \lr{epoch} خام بی‌معنا بود: مقدار $x^2$ به حدود $3.2\times10^{18}$ می‌رسد که از دقت ۵۳ بیتی اعداد اعشاری فراتر است و معادلات نرمال به صفر تقلیل می‌یابند. هر دو اصلاح شدند؛ دومی خطرناک‌تر بود چون عددی ظاهراً معقول تولید می‌کرد.

	\item \textbf{تداخل \lr{watchdog} با آزمایش.} در آزمایش ۲--۱ حالت «بیکار» با ۲۷ درصد مصرف پردازنده ثبت شد، چون \lr{watchdog} آشکارسازِ عمداً متوقف‌شده را ری‌استارت می‌کرد. با تعلیق موقت آستانه‌ی \lr{watchdog} فقط در همان بازه حل شد و مصرف پردازنده به ۱٫۳ درصد رسید.

	\item \textbf{منبع تغذیه‌ی ناکافی.} در بخش جداگانه‌ای توضیح داده شد.
\end{enumerate}

نکته‌ی روش‌شناختی که از این فهرست برمی‌آید: چند مورد از این ایرادها -- به‌ویژه شیب حداقل مربعات و باگ \lr{pipefail} -- نتایجی تولید می‌کردند که \emph{معقول به نظر می‌رسیدند}. یک عدد باورپذیر از یک محاسبه‌ی خراب، خطرناک‌تر از یک خطای آشکار است، و به همین دلیل در اسکریپت‌های آزمایش هرجا ممکن بود یک بررسی شاهد گنجانده شد: در آزمایش ۳--۶ اتصال کلاینت مجاز، در ۳--۷ ورود موفق با کلید، و در ۲--۱ بررسی اینکه حالت بیکار واقعاً بیکار باشد.
